\documentclass[conference]{IEEEtran}
\usepackage[T1]{fontenc}
\usepackage{amsmath}
\usepackage{booktabs}
\usepackage{array}
\usepackage{xcolor}
\usepackage[hidelinks]{hyperref}
\newcommand{\todo}[1]{\textcolor{red}{[#1]}}

\begin{document}

\title{Plausible but Wrong: What Verification Caught\\in a Small Adaptive-Mesh Research Code}

\author{\IEEEauthorblockN{\todo{Author name}}
\IEEEauthorblockA{\todo{Department}, University of California, Irvine\\
Irvine, CA, USA\\
\todo{email}}}

\maketitle

\begin{abstract}
We report on verifying a small research code (about 600 lines of Python)
written to compare three rules for choosing which cells an adaptive mesh
refines. Before any scientific claim was made, the code and the experimental
design contained eleven problems: two implementation defects, two
metric-definition defects, five design confounds and two tooling errors.
Every problem except the two tooling errors produced output that looked
plausible, and the five confounds each changed a conclusion. The defects were exposed by invariant
checks on mesh operations, convergence-order tests, and tracing every
reported number to a result file; the confounds were exposed by controlled
re-runs. The question the study set out to answer turned out to be the
smallest effect in the problem. We describe each problem, how it was
detected, and its measured effect, and draw lessons for testing research
software whose outputs have no exact oracle.
\end{abstract}

\begin{IEEEkeywords}
research software, scientific software testing, verification, adaptive mesh refinement, experience report
\end{IEEEkeywords}

\section{Introduction}
Research software errors are hard to see because their outputs are numbers
with no exact expected value, and a wrong number can look as reasonable as a
right one~\cite{hatton1994,kanewala2014}. Published results have been
retracted because of such errors~\cite{miller2006}. This paper is an
experience report from one small study in which we found and measured
eleven such problems before publication.

The study~\cite{amrpreprint} compared three \emph{marking rules} for
adaptive mesh refinement (AMR): a relative threshold, deterministic top-$K$,
and Gumbel top-$K$, a stochastic rule borrowed from machine
learning~\cite{kool2019}. The code solved the 1D viscous Burgers equation on
a periodic domain with a finite-volume scheme~\cite{leveque2002}, splitting
cells where an error indicator was large and merging them back where it was
small~\cite{berger1984}. The original question was which rule gives the
lowest error at a given budget.

Our contribution is a sourced log of each problem: what was wrong, whether
the output looked plausible, what detected it, and its measured effect
(Table~\ref{tab:log}). The code, the original scripts as received, the
superseded results and a script that re-measures each defect are public~\cite{repo}.

\section{The Code and How It Was Checked}
The code (\texttt{amr\_nonuniform.py}) stores a 1D mesh as sorted cell
edges. Each cell holds an average $u_i$ updated by flux differences,
$u_i \leftarrow u_i - \tfrac{\Delta t}{h_i}(F_{i+1/2}-F_{i-1/2})$, so mass
is conserved by construction. Refinement splits a cell in two; coarsening
merges two cells with a width-weighted average. A single global time step,
set by the smallest cell, advances every cell. We received three scripts
from a collaborator and kept them unchanged in \texttt{original/}; all
changes were made to copies.

Checks were added in three forms. \emph{Invariant checks}
(\texttt{check\_mesh.py}) assert properties that must hold after every
regrid: top-$K$ refines exactly $\min(K,\text{eligible})$ cells, no cell is
coarser than the starting mesh, every cell lies on the dyadic lattice, and
mass is conserved to round-off. \emph{Convergence tests}
(\texttt{check\_scheme.py}) confirm observed orders of accuracy near 1 and 2
on uniform grids and on a mesh with a refined patch, a standard code
verification practice~\cite{roache1998,oberkampf2010}. A \emph{number audit}
traces every number in the paper to the result file and column it came from.
Superseded outputs were kept, not deleted, so before-and-after values could
be compared row by row.

\section{What Was Found}

\begin{table*}[t]
\caption{Problems found, with measured effects. Plausible: did the output
look reasonable without the check? Type: D implementation defect, M metric
definition, C design confound, T tooling.}
\label{tab:log}
\centering
\small
\begin{tabular}{@{}l l p{4.3cm} c p{3.1cm} p{4.6cm}@{}}
\toprule
ID & Type & What was wrong & Plausible & Detected by & Measured effect \\
\midrule
D1 & D & Top-$K$ chose from all cells, then dropped those at the depth cap & yes & invariant: refinements $=\min(K,\text{eligible})$ & $K{=}25$: 25,25,25,21,20,21,21,20 per event (fixed: 25 each) \\
D2 & D & Coarsening merged any adjacent pair, not siblings & yes & invariants: width $\le$ base; dyadic mesh & largest cell $1.5\times$ base; 10.6\% of cells non-dyadic (fixed: $1.0\times$, 0\%) \\
M1 & M & Turnover averaged over empty$\to$empty and first-appearance pairs & yes & reviewing the definition & changed in all 154 adaptive runs per scheme; threshold mean 0.21$\to$0.15 \\
M2 & M & Budget search allowed refine threshold below coarsen threshold & no visible error & parameter audit & one configuration affected \\
C1 & C & First-order scheme's numerical viscosity ($\approx$0.02) exceeded physical $\nu{=}0.01$ & yes & analysis; second-order rerun & scheme changed error 5.7--7.0$\times$ vs.\ 3--30\% between rules \\
C2 & C & Rules compared at different final cell counts & yes & budget matching & residual mismatch (112 vs.\ 126 cells at target 128) \\
C3 & C & Cost measured by final cell count; global $\Delta t \propto h_{\min}^2$ & yes & work tracking; dense uniform sweep & 5-level wins at matched cells became losses at matched work \\
C4 & C & Refinement began after the front had started forming & yes & timing control & early refinement closed 61--82\% of the error gap \\
C5 & C & Gradient indicator refined smooth steep ramps & yes & indicator profiling & 103 of 126 refined cells in smooth regions \\
T1 & T & Syntax error in an experiment script & no (failed loudly) & run failed & none \\
T2 & T & Two experiments wrote the same output file & would overwrite silently & noticed before rerun & none \\
\bottomrule
\end{tabular}
\end{table*}

Table~\ref{tab:log} lists all eleven problems. Defect effects (D1, D2)
were measured by running the original and fixed code with identical
settings; the script is in the repository. We describe three in detail.

\subsection{D1: a budget parameter that did not set the budget}
Top-$K$ refines the $K$ highest-scoring cells, so $K$ is meant to control
how many cells are added. The original code chose the top $K$ among
\emph{all} cells and afterwards removed cells already at the maximum depth.
Once the front's cells reached that depth, fewer than $K$ cells were
refined per event (20--21 instead of 25 in our test). Nothing failed and
the solutions looked reasonable; the error was visible only to an invariant
comparing refinements per event with $\min(K,\text{eligible})$. Because the
study's comparisons depended on controlling budgets, this defect undermined
the experimental control rather than the numerics.

\subsection{C1 and C3: confounds that changed the answer}
Some of the most consequential problems were not bugs. The first-order
scheme adds numerical viscosity of about $\alpha h/2$, roughly 0.02 on the
base cells, twice the physical viscosity. Coarse regions therefore solved a
different equation, and at matched work a uniform grid beat every adaptive
configuration (21 of 21). With a second-order scheme the error fell by
5.7--7.0$\times$, while the three rules differed by only 3--30\%.

Cost was a second confound. Matching rules on final cell count looked fair,
but the global time step is limited by $h_{\min}^2/(2\nu)$ once cells are
small, so each extra level of refinement can multiply the number of steps
for \emph{every} cell by up to four. Final cell count also ignores how many
cells existed during the run. Comparing at equal work, measured as total
cell updates against uniform runs that bracket each adaptive run's work,
turned five-level configurations that won at matched cell count into
losses.

\subsection{M1: a metric defined by convention}
Grid turnover, the Jaccard distance between finest-level regions at
consecutive regrids, was averaged over all pairs. A pair of empty sets
contributed 0 and a first appearance contributed 1, so runs that never
reached the finest level reported a turnover of exactly 0. Restricting to
pairs where both sets are non-empty changed every turnover value; accuracy
values were unchanged, which we confirmed row by row against the superseded
files.

\section{Lessons}
\textbf{Test the experimental controls, not only the numerics.}
Convergence tests passed throughout; D1 and D2 were invisible to them. The
defects lived in the code that implemented the \emph{experiment}, the
budget and mesh bookkeeping. Invariants on those operations were cheap and
caught both. Property-based testing of mesh operations is a natural
extension.

\textbf{Make every claimed number traceable.} The number audit forced each
claim back to a file, which exposed M1 and kept superseded values out of
the paper. Keeping superseded outputs made the effect of each fix
measurable rather than asserted.

\textbf{Choose cost metrics that see global effects.} A local decision (one
small cell) had a global cost (every cell's time step). A metric that
counts only what is visible in the final state can reverse a comparison.

\textbf{Check the design before comparing methods.} Five of the eleven
problems were design confounds, and they mattered more than any bug: the
study's original question turned out to be the smallest effect.

\textbf{AI assistance.} \todo{Describe accurately: which checks the AI
assistant proposed, which errors it introduced (e.g., T1, T2 if accurate),
and what the author verified by hand. Keep only what you can support.}

\section{Threats to Validity}
This is a single case study of one small code, written up after the fact by
the people who found the problems, with no comparison group. The order in
which problems were discovered is not recorded in version control, which
begins after the fixes, so we make no claims about timing. The effect of D2
on accuracy was not isolated. Four earlier defects described in the
original author's code comments were fixed before we received the code; we
did not verify them and do not count them. The lessons are offered as
hypotheses for research codes of similar size, not as general findings.

\section*{Acknowledgment}
\todo{Collaborator who supplied the original scripts, if they agree.}
\todo{AI-use disclosure following the ICSE 2027 policy.}

\begin{thebibliography}{10}
\bibitem{hatton1994}
L.~Hatton and A.~Roberts, ``How accurate is scientific software?''
\emph{IEEE Trans. Softw. Eng.}, vol.~20, no.~10, pp.~785--797, 1994.

\bibitem{kanewala2014}
U.~Kanewala and J.~M. Bieman, ``Testing scientific software: A systematic
literature review,'' \emph{Inf. Softw. Technol.}, vol.~56, no.~10,
pp.~1219--1232, 2014.

\bibitem{miller2006}
G.~Miller, ``A scientist's nightmare: Software problem leads to five
retractions,'' \emph{Science}, vol.~314, no.~5807, pp.~1856--1857, 2006.

\bibitem{amrpreprint}
\todo{Author}, ``AMR marking strategies for viscous Burgers,'' arXiv
preprint \todo{arXiv ID}, 2026.

\bibitem{kool2019}
W.~Kool, H.~van Hoof, and M.~Welling, ``Stochastic beams and where to find
them: The Gumbel-top-$k$ trick for sampling sequences without replacement,''
in \emph{Proc. ICML}, 2019, pp.~3499--3508.

\bibitem{leveque2002}
R.~J. LeVeque, \emph{Finite Volume Methods for Hyperbolic Problems}.
Cambridge University Press, 2002.

\bibitem{berger1984}
M.~J. Berger and J.~Oliger, ``Adaptive mesh refinement for hyperbolic
partial differential equations,'' \emph{J. Comput. Phys.}, vol.~53, no.~3,
pp.~484--512, 1984.

\bibitem{repo}
\todo{Author}, ``amr-marking-study,'' GitHub repository,
\url{https://github.com/achyutarag/amr-marking-study}, 2026,
\todo{Zenodo DOI}.

\bibitem{roache1998}
P.~J. Roache, \emph{Verification and Validation in Computational Science
and Engineering}. Albuquerque, NM: Hermosa, 1998.

\bibitem{oberkampf2010}
W.~L. Oberkampf and C.~J. Roy, \emph{Verification and Validation in
Scientific Computing}. Cambridge University Press, 2010.
\end{thebibliography}

\end{document}
